%%%%%%%%%%%%%%%%%%%%%%%%%%%%%%%%%%%%%%%%%%%%%%%%%%%%%%%%%%%%%%%%%%%%%%%%%%%%%%%
% SECTION 11.11: Psychological Distortions and Group Effects
% Axioms: AWX-A66 to AWX-A70a
% Appendix UNMAPPED_AWA Reference: Section 3.9
%%%%%%%%%%%%%%%%%%%%%%%%%%%%%%%%%%%%%%%%%%%%%%%%%%%%%%%%%%%%%%%%%%%%%%%%%%%%%%%

Awareness can be systematically distorted by psychological defense mechanisms 
and group dynamics. These distortions are not random errors but structured 
patterns that serve psychological functions—often at the cost of accurate 
perception of consequences.

\subsubsection{Individual Defense Mechanisms (AWX-A66 to A69)}

\paragraph{AWX-A66: Suppression (Verdrängung).}
Awareness of threatening consequences can be actively suppressed:
\begin{equation}
A_{suppressed} = A_{base} \cdot (1 - \sigma \cdot Threat)
\end{equation}
where $\sigma$ is suppression strength and $Threat$ measures how threatening 
the consequence is to self-image or wellbeing.

Key characteristics:
\begin{itemize}
    \item Suppression is \emph{motivated}: more threatening = more suppressed
    \item Suppression is effortful: requires cognitive resources
    \item Suppression is unstable: suppressed awareness can resurface
    \item Suppression correlates with $A_{IDN}$: identity-threatening 
          information is most suppressed
\end{itemize}

\paragraph{AWX-A67: Splitting (Spaltung).}
Awareness can be compartmentalized, preventing integration across domains:
\begin{equation}
A_{domain_1} \perp A_{domain_2}
\end{equation}

This explains how someone can be simultaneously aware of climate change 
(in abstract discussion) and unaware of it (when booking a flight). 
The awareness exists but is not integrated.

\paragraph{AWX-A68: Norm Conflict.}
When multiple norms apply and conflict, awareness becomes paralyzed:
\begin{equation}
A_{effective} = A_{base} \cdot \frac{1}{1 + |Norm_1 - Norm_2|}
\end{equation}

When norms strongly conflict, effective awareness approaches zero, 
leading to inaction or random choice.

\paragraph{AWX-A69: False Awareness.}
Attention can be misdirected to irrelevant consequences while missing 
relevant ones:
\begin{equation}
A_{irrelevant} \gg A_{relevant}
\end{equation}

This is particularly problematic when salient but unimportant factors 
capture attention (e.g., brand logos) while important but less salient 
factors (e.g., terms and conditions) are ignored.

\subsubsection{Group-Level Distortions (AWX-A70 to A70a)}

\paragraph{AWX-A70: Collective Override (Kollektive Überlagerung).}
In group settings, collective awareness can override individual awareness:
\begin{equation}
A_{effective} = w_{group} \cdot A_{group} + (1 - w_{group}) \cdot A_{individual}
\end{equation}

When group identification is strong ($w_{group} \to 1$), individual 
awareness becomes irrelevant. This explains:
\begin{itemize}
    \item Groupthink in organizational decisions
    \item Mob behavior in crowds
    \item Conformity in social settings
    \item Tribal dynamics in political contexts
\end{itemize}

\paragraph{AWX-A70a: Mirror Gap (Spiegelkohärenz).}
A systematic discrepancy between self-awareness and perceived awareness:
\begin{equation}
A^{self} \neq A^{perceived}
\end{equation}

Specifically:
\begin{itemize}
    \item We overestimate others' awareness of our internal states 
          (spotlight effect)
    \item We underestimate the degree to which others are aware of 
          their own consequences (false uniqueness)
    \item We assume others' awareness patterns mirror our own 
          (false consensus)
\end{itemize}

\subsubsection{Distortion Interactions}

The distortion mechanisms interact in complex ways:

\begin{table}[htbp]
\centering
\caption{Distortion Interaction Patterns}
\label{tab:distortion-interactions}
\begin{tabular}{@{}lll@{}}
\toprule
\textbf{Combination} & \textbf{Effect} & \textbf{Example} \\
\midrule
Suppression + Splitting & Compartmentalized denial & Addiction awareness \\
Norm conflict + Group & Polarization paralysis & Political tribalism \\
False awareness + Mirror & Mutual misunderstanding & Organizational silos \\
Suppression + Group & Collective denial & Corporate misconduct \\
\bottomrule
\end{tabular}
\end{table}

\subsubsection{Clinical and Subclinical Implications}

While the AWX framework describes normal psychological functioning, 
the distortion mechanisms have clinical analogs:

\begin{itemize}
    \item \textbf{AWX-A66} (Suppression) relates to repression and denial
    \item \textbf{AWX-A67} (Splitting) relates to dissociation
    \item \textbf{AWX-A68} (Norm conflict) relates to cognitive dissonance
    \item \textbf{AWX-A69} (False awareness) relates to attention disorders
    \item \textbf{AWX-A70} (Collective override) relates to deindividuation
\end{itemize}

The BCM framework does not pathologize these mechanisms—they are normal 
and often adaptive. However, extreme manifestations can become dysfunctional.


\paragraph{AWX-A10b: Clinical Relevance Threshold.}
AWX-A10b provides a formal criterion for when distortion patterns become 
clinically significant:

\begin{axiom}[AWX-A10b: Klinische Relevanz]
Persistent hyper-awareness or blockade states that exceed a time threshold 
and cause distress are clinically relevant.

\textbf{Formal Condition:}
\begin{equation}
\text{Clinical Relevance} \Leftrightarrow 
\text{Persistence} > \theta_{time} \land \text{Distress} > \theta_{distress}
\end{equation}
\end{axiom}

\paragraph{Application.}
This axiom distinguishes:
\begin{itemize}
    \item \textbf{Normal variation}: Temporary suppression or heightened 
          awareness within adaptive bounds
    \item \textbf{Subclinical patterns}: Persistent but low-distress 
          deviations (e.g., chronic mild avoidance)
    \item \textbf{Clinical conditions}: Persistent high-distress patterns 
          requiring professional intervention
\end{itemize}

The threshold parameters ($\theta_{time}$, $\theta_{distress}$) are context-
dependent and must be calibrated against clinical standards. BCM does not 
diagnose—it provides a framework for understanding when awareness patterns 
deviate from adaptive functioning.

\subsubsection{Implications for Intervention}

Understanding distortions suggests intervention strategies:

\begin{enumerate}
    \item \textbf{Reduce threat}: Lower perceived threat to reduce 
          suppression motivation
    \item \textbf{Bridge compartments}: Create connections between 
          split awareness domains
    \item \textbf{Resolve norm conflicts}: Clarify which norms apply 
          in specific contexts
    \item \textbf{Redirect attention}: Make relevant consequences 
          more salient than irrelevant ones
    \item \textbf{Individuate}: Reduce group pressure to allow 
          individual awareness to surface
    \item \textbf{Calibrate mirrors}: Provide feedback on actual 
          (vs. perceived) awareness of others
\end{enumerate}

\subsubsection{Measurement Challenges}

Distortions create measurement challenges:
\begin{itemize}
    \item Self-reported awareness may reflect \emph{suppressed} awareness
    \item Behavioral measures may reflect \emph{split} awareness 
          (aware in survey, unaware in action)
    \item Group measures may reflect \emph{collective} rather than 
          individual awareness
    \item Awareness measures at one moment may not predict awareness 
          at the Moment of Truth
\end{itemize}

This is why CASE-09 (Sustainable Consumption) shows a 4.7:1 ratio between 
survey awareness and store awareness—the measurement context triggers 
different distortion patterns.
